\documentclass[12pt, a4paper]{article}

% ==========================================
% 1. COMPILER, LANGUAGE & FONT SETUP
% ==========================================
\usepackage{fontspec}
\usepackage{polyglossia}
\setmainlanguage{bengali}
\setotherlanguage{english}
\newfontfamily\bengalifont[Script=Bengali, AutoFakeBold=2.5, AutoFakeSlant=0.2]{Kalpurush}

% ==========================================
% 2. MATHEMATICS, UNITS & FORMATTING
% ==========================================
\usepackage{amsmath, amssymb, siunitx}
\usepackage[margin=1in]{geometry}
\usepackage{setspace}
\usepackage{enumitem}
\usepackage{microtype} % For better typography
\onehalfspacing % Better line spacing for readability

% ==========================================
% 3. COLORS & BEAUTIFICATION PACKAGES
% ==========================================
\usepackage[dvipsnames, table]{xcolor}
\usepackage[skins, breakable, theorems]{tcolorbox}
\usepackage{tikz}
\renewcommand{\labelitemi}{$\bullet$}

% Define custom beautiful colors
\definecolor{primary}{HTML}{0056b3}
\definecolor{secondary}{HTML}{e7f1ff}
\definecolor{success}{HTML}{198754}
\definecolor{successbg}{HTML}{d1e7dd}
\definecolor{accent}{HTML}{fd7e14}
\definecolor{darkgray}{HTML}{343a40}

% Custom Tcolorbox Styles
\tcbset{
    questionbox/.style={
        enhanced, colback=secondary, colframe=primary, arc=2mm, boxrule=1pt,
        fonttitle=\bfseries\Large, title=#1,
        drop fuzzy shadow=black!10,
        attach boxed title to top left={xshift=5mm, yshift=-3mm},
        boxed title style={colback=primary, arc=1mm}
    },
    answerbox/.style={
        enhanced, colback=successbg, colframe=success, arc=1.5mm, boxrule=1pt,
        left=2mm, right=2mm, top=2mm, bottom=2mm,
        drop fuzzy shadow=black!10
    },
    solutionbox/.style={
        enhanced, colback=white, colframe=darkgray, arc=2mm, boxrule=1pt,
        fonttitle=\bfseries\large, title=#1, breakable,
        coltitle=white, colbacktitle=darkgray,
        drop fuzzy shadow=black!10,
        attach boxed title to top center={yshift=-3mm},
        boxed title style={arc=1mm}
    }
}

\begin{document}

\newpage
% ------------------------------------------
% QUESTION SECTION 1 - DETAILED & CLASSY
% ------------------------------------------
\begin{tcolorbox}[questionbox={প্রশ্ন (Question) 1}]
\noindent
একটি অগ্রগামী তরঙ্গের সমীকরণ $y_1 = 15 \sin \frac{2\pi}{6}(100t - x)\text{ m}$। তরঙ্গটি মুক্ত প্রান্তে প্রতিফলিত হয়ে বিপরীত দিকে ফিরে এলে লব্ধি স্থির তরঙ্গের সমীকরণ গঠন করো এবং $x = 3\text{ m}$ অবস্থানে অবস্থিত কণার পর্যায়কালের বিভিন্ন মুহূর্তে ($t = 0, \frac{T}{4}, \frac{T}{2}, \frac{3T}{4}, T$) সরণের মান নির্ণয় করো। এছাড়াও উক্ত কণার জন্য সরণ-সময় ($y\text{ vs }t$) এবং বেগ-সময় ($v\text{ vs }t$) লেখচিত্র অঙ্কন করো।

\vspace{0.8em}
\begin{english}
\noindent\textit{\color{darkgray}
The equation of a progressive wave is $y_1 = 15 \sin \frac{2\pi}{6}(100t - x)\text{ m}$. When the wave reflects from a free boundary and returns in the opposite direction, form the equation of the resulting standing wave, and calculate the displacement of a particle at $x = 3\text{ m}$ at different fractions of time period $t = 0, \frac{T}{4}, \frac{T}{2}, \frac{3T}{4}, T$. Also draw the displacement-time ($y\text{ vs }t$) and velocity-time ($v\text{ vs }t$) graphs for the particle.
}
\end{english}
\end{tcolorbox}

\vspace{1em}

% ------------------------------------------
% SOLUTION SECTION 1
% ------------------------------------------
\begin{tcolorbox}[solutionbox={সমাধান (Solution)}]
\noindent
১. স্থির তরঙ্গের সমীকরণ গঠন: মুক্ত প্রান্তে প্রতিফলনের ফলে কোনো দশা পরিবর্তন ($\Delta \phi = 0$) হয় না। প্রতিফলিত তরঙ্গ: $y_2 = 15 \sin \frac{2\pi}{6}(100t + x)$ \\
লব্ধি স্থির তরঙ্গ:
\[ y = y_1 + y_2 = 15 \left[ \sin \frac{2\pi}{6}(100t - x) + \sin \frac{2\pi}{6}(100t + x) \right] \]
\[ y = 30 \cos\left(\frac{2\pi x}{6}\right) \sin\left(\frac{2\pi \times 100t}{6}\right) = 30 \cos\left(\frac{\pi x}{3}\right) \sin\left(\frac{100\pi t}{3}\right)\text{ m} \]

\noindent
২. $x = 3\text{ m}$ কণার জন্য সরণ ও বেগ সমীকরণ:
\[ A(x=3) = 30 \cos\left(\frac{3\pi}{3}\right) = 30 \cos(\pi) = -30\text{ m} \]
সুতরাং, সরণের সমীকরণ: 
\[ y(3, t) = -30 \sin\left(\frac{100\pi t}{3}\right) = -30 \sin\left(\frac{2\pi t}{T}\right)\text{ m} \]
কণাটির বেগের সমীকরণ ($v = \frac{dy}{dt}$):
\[ v(3, t) = -30 \times \frac{100\pi}{3} \cos\left(\frac{100\pi t}{3}\right) = -1000\pi \cos\left(\frac{2\pi t}{T}\right)\text{ m s}^{-1} \]

\noindent
৩. বিভিন্ন সময়ে সরণ ($y$) ও বেগ ($v$) এর মান: \\
- $t = 0$ হলে: $y = 0\text{ m}$, \quad $v = -1000\pi \approx -3141.59\text{ m s}^{-1}$ \\
- $t = \frac{T}{4}$ হলে: $y = -30\text{ m}$, \quad $v = 0\text{ m s}^{-1}$ \\
- $t = \frac{T}{2}$ হলে: $y = 0\text{ m}$, \quad $v = +1000\pi \approx +3141.59\text{ m s}^{-1}$ \\
- $t = \frac{3T}{4}$ হলে: $y = +30\text{ m}$, \quad $v = 0\text{ m s}^{-1}$ \\
- $t = T$ হলে: $y = 0\text{ m}$, \quad $v = -1000\pi \approx -3141.59\text{ m s}^{-1}$

\vspace{1em}
\noindent
\textbf{৪. সরণ-সময় ($y\text{ vs }t$) ও বেগ-সময় ($v\text{ vs }t$) লেখচিত্র:}

\begin{center}
\begin{tikzpicture}[scale=1.05, >=stealth]
    % ==========================================
    % GRAPH 1: Displacement vs Time (y vs t)
    % ==========================================
    \begin{scope}[shift={(0,0)}]
        % Grid
        \draw[gray!20, thin] (-0.5,-2.2) grid (5.5,2.2);

        % Axes
        \draw[->, thick] (-0.5,0) -- (5.8,0) node[right] {$t$};
        \draw[->, thick] (0,-2.3) -- (0,2.3) node[above] {$y\text{ (m)}$};

        % Y-Axis Ticks & Labels
        \draw (0.1, 1.5) -- (-0.1, 1.5) node[left, font=\small] {$+30$};
        \draw (0.1, -1.5) -- (-0.1, -1.5) node[left, font=\small] {$-30$};
        \node[below left, font=\small] at (0,0) {$0$};

        % X-Axis Ticks & Labels
        \draw (1.25, 0.1) -- (1.25, -0.1) node[above=2pt, font=\small] {$\frac{T}{4}$};
        \draw (2.5, 0.1) -- (2.5, -0.1) node[below=2pt, font=\small] {$\frac{T}{2}$};
        \draw (3.75, 0.1) -- (3.75, -0.1) node[below=2pt, font=\small] {$\frac{3T}{4}$};
        \draw (5.0, 0.1) -- (5.0, -0.1) node[below=2pt, font=\small] {$T$};

        % Dashed Projection Lines
        \draw[dashed, red!60] (0, -1.5) -- (1.25, -1.5) -- (1.25, 0);
        \draw[dashed, blue!60] (0, 1.5) -- (3.75, 1.5) -- (3.75, 0);

        % Curve: y = -1.5 * sin(2*pi*x / 5)
        \draw[ultra thick, blue!80!black, domain=0:5, samples=100] 
            plot (\x, {-1.5*sin(\x*360/5)});

        % Key Points
        \fill[red!80!black] (0,0) circle (2.5pt);
        \fill[red!80!black] (1.25,-1.5) circle (2.5pt);
        \fill[red!80!black] (2.5,0) circle (2.5pt);
        \fill[red!80!black] (3.75,1.5) circle (2.5pt);
        \fill[red!80!black] (5.0,0) circle (2.5pt);

        % Graph Title
        \node[above, blue!80!black, font=\small\bfseries] at (2.5, 2.3) {সরণ-সময় লেখচিত্র ($y\text{ vs }t$)};
    \end{scope}

    % ==========================================
    % GRAPH 2: Velocity vs Time (v vs t)
    % ==========================================
    \begin{scope}[shift={(0,-5.5)}]
        % Grid
        \draw[gray!20, thin] (-0.5,-2.2) grid (5.5,2.2);

        % Axes
        \draw[->, thick] (-0.5,0) -- (5.8,0) node[right] {$t$};
        \draw[->, thick] (0,-2.3) -- (0,2.3) node[above] {$v\text{ (m s}^{-1})$};

        % Y-Axis Ticks & Labels
        \draw (0.1, 1.5) -- (-0.1, 1.5) node[left, font=\small] {$+1000\pi$};
        \draw (0.1, -1.5) -- (-0.1, -1.5) node[left, font=\small] {$-1000\pi$};
        \node[below left, font=\small] at (0,0) {$0$};

        % X-Axis Ticks & Labels
        \draw (1.25, 0.1) -- (1.25, -0.1) node[below=2pt, font=\small] {$\frac{T}{4}$};
        \draw (2.5, 0.1) -- (2.5, -0.1) node[below=2pt, font=\small] {$\frac{T}{2}$};
        \draw (3.75, 0.1) -- (3.75, -0.1) node[below=2pt, font=\small] {$\frac{3T}{4}$};
        \draw (5.0, 0.1) -- (5.0, -0.1) node[below=2pt, font=\small] {$T$};

        % Dashed Projection Lines
        \draw[dashed, red!60] (0, -1.5) -- (5.0, -1.5);
        \draw[dashed, blue!60] (0, 1.5) -- (2.5, 1.5) -- (2.5, 0);

        % Curve: v = -1.5 * cos(2*pi*x / 5)
        \draw[ultra thick, teal!80!black, domain=0:5, samples=100] 
            plot (\x, {-1.5*cos(\x*360/5)});

        % Key Points
        \fill[red!80!black] (0,-1.5) circle (2.5pt);
        \fill[red!80!black] (1.25,0) circle (2.5pt);
        \fill[red!80!black] (2.5,1.5) circle (2.5pt);
        \fill[red!80!black] (3.75,0) circle (2.5pt);
        \fill[red!80!black] (5.0,-1.5) circle (2.5pt);

        % Graph Title
        \node[above, teal!80!black, font=\small\bfseries] at (2.5, 2.3) {বেগ-সময় লেখচিত্র ($v\text{ vs }t$)};
    \end{scope}
\end{tikzpicture}
\end{center}
\end{tcolorbox}

% ------------------------------------------
% QUESTION SECTION 12 - DETAILED & CLASSY
% ------------------------------------------
\begin{tcolorbox}[questionbox={প্রশ্ন (Question) 2}]
\noindent
একটি রকেট ভূমির সাথে উল্লম্বভাবে $200\text{ km h}^{-1}$ সুষম বেগে ওপরের দিকে উঠছে। ভূমির ওপর নির্দিষ্ট উচ্চতায় পৌঁছানোর মুহূর্তে রকেট থেকে ভূমির দিকে একটি সংকেত (শব্দ) পাঠানো হলো। ভূমি থেকে প্রতিফলিত প্রতিধ্বনি রকেটে পৌঁছাতে $3.3\text{ s}$ সময় নেয়। বাতাসে শব্দের বেগ $340\text{ m s}^{-1}$ হলে, সংকেত পাঠানোর মুহূর্তে রকেটের উচ্চতা কত ছিল?

\vspace{0.8em}
\begin{english}
\noindent\textit{\color{darkgray}
A rocket ascends vertically upwards at a uniform speed of $200\text{ km h}^{-1}$. At a certain height above the ground, a sound signal is emitted from the rocket towards the ground. The echo reflected from the ground reaches the rocket after $3.3\text{ s}$. If the speed of sound in air is $340\text{ m s}^{-1}$, what was the altitude of the rocket at the instant the signal was emitted?
}
\end{english}
\end{tcolorbox}

\vspace{1em}
% ------------------------------------------
% SOLUTION SECTION 12
% ------------------------------------------
\begin{tcolorbox}[solutionbox={সমাধান (Solution)}]
\noindent
১. প্রদত্ত উপাত্ত: \\
- রকেটের বেগ, $v_r = \frac{200}{3.6} = 55.56\text{ m s}^{-1}$ \\
- শব্দের বেগ, $v_s = 340\text{ m s}^{-1}$ \\
- মোট সময়, $t = 3.3\text{ s}$ \\
২. দূরত্ব সমীকরণ গঠন: ধরি, সংকেত নিসরণের মুহূর্তে রকেটের উচ্চতা ছিল $h$। $t = 3.3\text{ s}$ সময়ে রকেটটি ওপরের দিকে অতিক্রান্ত দূরত্ব, $s_r = v_r \times t = 55.56 \times 3.3 = 183.35\text{ m}$। \\
শব্দ তরঙ্গটিকে নিচে ভূমিতে গিয়ে ($h$) আবার প্রতিফলিত হয়ে রকেটের নতুন অবস্থানে পৌঁছাতে অতিক্রান্ত দূরত্ব = $h + (h - s_r) = 2h - s_r$। \\
শব্দের অতিক্রান্ত মোট দূরত্ব:
\[ 2h - s_r = v_s \times t \]
\[ 2h - 183.35 = 340 \times 3.3 = 1122\text{ m} \]
\[ 2h = 1122 + 183.35 = 1305.35\text{ m} \]
\[ h = \frac{1305.35}{2} = 652.68\text{ m} \]
অতএব, সংকেত পাঠানোর মুহূর্তে রকেটের উচ্চতা ছিল $652.68\text{ m}$।
\end{tcolorbox}

\vspace{1.5em}

% ------------------------------------------
% QUESTION SECTION 13 - DETAILED & CLASSY
% ------------------------------------------
\begin{tcolorbox}[questionbox={প্রশ্ন (Question) 3}]
\noindent
$1\text{ m}$ দীর্ঘ একটি তার $40\text{ N}$ বল দ্বারা টানা অবস্থায় $320\text{ Hz}$ কম্পাঙ্কের একটি সুরশলাকার সাথে ঐকতানে (Resonance) কাঁপছে। তারের টান পরিবর্তন করে $20\text{ N}$ করা হলে, তারটিকে পুনরায় উক্ত সুরশলাকার সাথে ঐকতানে আনতে এর দৈর্ঘ্য কতটুকু পরিবর্তন করতে হবে?

\vspace{0.8em}
\begin{english}
\noindent\textit{\color{darkgray}
A wire of length $1\text{ m}$ under a tension of $40\text{ N}$ vibrates in resonance with a tuning fork of frequency $320\text{ Hz}$. If the tension is changed to $20\text{ N}$, by how much must the length of the wire be altered to bring it back into resonance with the same tuning fork?
}
\end{english}
\end{tcolorbox}

\vspace{1em}

% ------------------------------------------
% SOLUTION SECTION 13
% ------------------------------------------
\begin{tcolorbox}[solutionbox={সমাধান (Solution)}]
\noindent
১. সূত্র প্রয়োগ: টানা তারের মূলসুরের কম্পাঙ্কের সূত্র: $f = \frac{1}{2l}\sqrt{\frac{T}{\mu}}$। যেহেতু কম্পাঙ্ক $f = 320\text{ Hz}$ এবং একক দৈর্ঘ্যের ভর $\mu$ অপরিবর্তিত থাকে:
\[ \frac{\sqrt{T_1}}{l_1} = \frac{\sqrt{T_2}}{l_2} \]
২. নতুন দৈর্ঘ্য ($l_2$) নির্ণয়: প্রদত্ত, $T_1 = 40\text{ N}$, $l_1 = 1.0\text{ m}$, $T_2 = 20\text{ N}$।
\[ l_2 = l_1 \times \sqrt{\frac{T_2}{T_1}} = 1.0 \times \sqrt{\frac{20}{40}} = \frac{1}{\sqrt{2}} \approx 0.7071\text{ m} = 70.71\text{ cm} \]
৩. দৈর্ঘ্যের পরিবর্তন ($\Delta l$) নির্ণয়:
\[ \Delta l = l_1 - l_2 = 1.0 - 0.7071 = 0.2929\text{ m} = 29.29\text{ cm} \]
অতএব, তারটিকে পুনরায় ঐকতানে আনতে এর দৈর্ঘ্য $29.29\text{ cm}$ কমাতে হবে।
\end{tcolorbox}

\vspace{1.5em}

% ------------------------------------------
% QUESTION SECTION 14 - DETAILED & CLASSY
% ------------------------------------------
\begin{tcolorbox}[questionbox={প্রশ্ন (Question) 4}]
\noindent
$A, B, C$ ও $D$ চারটি সুরশলাকা রয়েছে। $A$ শলাকার কম্পাঙ্ক $250\text{ Hz}$। $A$ শলাকাটি $B$ ও $D$-এর সাথে যথাক্রমে প্রতি সেকেন্ডে $2$ টি ও $6$ টি বিট উৎপন্ন করে। $B$ ও $D$ পরস্পরের সাথে প্রতি সেকেন্ডে $4$ টি বিট উৎপন্ন করে। আবার, $B$ ও $D$ শলাকাদ্বয়ের প্রতিটি $C$-এর সাথে প্রতি সেকেন্ডে $2$ টি বিট উৎপন্ন করে। $C$ সুরশলাকার সম্ভাব্য কম্পাঙ্ক গাণিতিকভাবে বিশ্লেষণ করো।

\vspace{0.8em}
\begin{english}
\noindent\textit{\color{darkgray}
Four tuning forks $A, B, C,$ and $D$ are given. The frequency of $A$ is $250\text{ Hz}$. Fork $A$ produces $2$ and $6$ beats per second with $B$ and $D$ respectively. $B$ and $D$ produce $4$ beats per second with each other. Moreover, each of the forks $B$ and $D$ produces $2$ beats per second with $C$. Analyze mathematically to find the possible frequencies of tuning fork $C$.
}
\end{english}
\end{tcolorbox}

\vspace{1em}

% ------------------------------------------
% SOLUTION SECTION 14
% ------------------------------------------
\begin{tcolorbox}[solutionbox={সমাধান (Solution)}]
\noindent
১. $B$ ও $D$-এর কম্পাঙ্ক নির্ধারণ: দেওয়া আছে, $f_A = 250\text{ Hz}$। \\
- $f_B = 250 \pm 2 = 252\text{ Hz}$ অথবা $248\text{ Hz}$ \\
- $f_D = 250 \pm 6 = 256\text{ Hz}$ অথবা $244\text{ Hz}$ \\
যেহেতু $B$ ও $D$ প্রতি সেকেন্ডে $4$ টি বিট দেয় ($|f_B - f_D| = 4$), তাই দুটি ক্ষেত্র সম্ভব: \\
- ক্ষেত্র-১: $f_B = 252\text{ Hz}$ এবং $f_D = 256\text{ Hz}$ \\
- ক্ষেত্র-২: $f_B = 248\text{ Hz}$ এবং $f_D = 244\text{ Hz}$ \\
২. $C$-এর সম্ভাব্য কম্পাঙ্ক গাণিতিক পরীক্ষা: $C$ শলাকাটি $B$ এবং $D$ উভয়ের সাথেই প্রতি সেকেন্ডে $2$ টি বিট সৃষ্টি করে। \\
- ক্ষেত্র-১ এর জন্য ($f_B = 252\text{ Hz}, f_D = 256\text{ Hz}$): \\
$|f_C - 252| = 2 \implies f_C = 254\text{ Hz}$ অথবা $250\text{ Hz}$। এখন $f_D = 256$-এর সাথে মেলালে: $|254 - 256| = 2$ (সঠিক)। (কিন্তু $f_C = 250$ হলে $|250 - 256| = 6 \neq 2$)। সুতরাং, ১ম ক্ষেত্রে $f_C = 254\text{ Hz}$। \\
- ক্ষেত্র-২ এর জন্য ($f_B = 248\text{ Hz}, f_D = 244\text{ Hz}$): \\
$|f_C - 248| = 2 \implies f_C = 246\text{ Hz}$ অথবা $250\text{ Hz}$। এখন $f_D = 244$-এর সাথে মেলালে: $|246 - 244| = 2$ (সঠিক)। সুতরাং, ২য় ক্ষেত্রে $f_C = 246\text{ Hz}$। \\
অতএব, $C$ সুরশলাকাটির সম্ভাব্য কম্পাঙ্ক হলো $254\text{ Hz}$ অথবা $246\text{ Hz}$।
\end{tcolorbox}

\vspace{1.5em}

% ------------------------------------------
% QUESTION SECTION 15 - DETAILED & CLASSY
% ------------------------------------------
\begin{tcolorbox}[questionbox={প্রশ্ন (Question) 5}]
\noindent
$256\text{ Hz}$ কম্পাঙ্কের একটি সুরশলাকা একটি দুই মুখ খোলা নলের সাথে প্রতি সেকেন্ডে $8$ টি বিট উৎপন্ন করে। সুরশলাকাটির বাহুতে মোম লাগালে বিট সংখ্যা বৃদ্ধি পায়। বাতাসে শব্দের বেগ $285.25\text{ m s}^{-1}$ হলে নলটির প্রাথমিক দৈর্ঘ্য কত? নলটিকে এক মুখ বন্ধ অর্গান নলে রূপান্তরিত করে সুরশলাকাটির (মোম ছাড়া) সাথে অনুনাদ সৃষ্টি করতে হলে নলের দৈর্ঘ্যের পরিবর্তন কত হবে?

\vspace{0.8em}
\begin{english}
\noindent\textit{\color{darkgray}
A tuning fork of frequency $256\text{ Hz}$ produces $8$ beats per second with an open organ pipe. Loading wax on the tuning fork increases the beat frequency. If the speed of sound in air is $285.25\text{ m s}^{-1}$, what is the initial length of the pipe? If the pipe is converted into a closed organ pipe to resonate with the unloaded tuning fork, what must be the change in length of the pipe?
}
\end{english}
\end{tcolorbox}

\vspace{1em}

% ------------------------------------------
% SOLUTION SECTION 15
% ------------------------------------------
\begin{tcolorbox}[solutionbox={সমাধান (Solution)}]
\noindent
১. খোলা নলের কম্পাঙ্ক ($f_{\text{pipe}}$) নির্ধারণ: সুরশলাকার কম্পাঙ্ক $f_{\text{fork}} = 256\text{ Hz}$। বিট সংখ্যা = $8 \implies f_{\text{pipe}} = 256 + 8 = 264\text{ Hz}$ অথবা $256 - 8 = 248\text{ Hz}$। \\
সুরশলাকায় মোম লাগালে এর কম্পাঙ্ক কমে। যেহেতু এর ফলে বিট সংখ্যা বৃদ্ধি পায়, তাই নলের কম্পাঙ্ক অবশ্যই সুরশলাকার কম্পাঙ্কের চেয়ে বেশি ছিল। সুতরাং, $f_{\text{pipe}} = 264\text{ Hz}$। \\
২. খোলা নলের প্রাথমিক দৈর্ঘ্য ($L_1$) নির্ণয়:
\[ f_{\text{pipe}} = \frac{v}{2 L_1} \implies 264 = \frac{285.25}{2 L_1} \]
\[ L_1 = \frac{285.25}{2 \times 264} = 0.5402\text{ m} = 54.02\text{ cm} \]
৩. এক মুখ বন্ধ নলে রূপান্তরিত করার পর নতুন দৈর্ঘ্য ($L_2$) নির্ণয়: সুরশলাকা ($256\text{ Hz}$)-এর সাথে বন্ধ নলের মূলসুরের অনুনাদের জন্য:
\[ f_{\text{closed}} = \frac{v}{4 L_2} \implies 256 = \frac{285.25}{4 L_2} \]
\[ L_2 = \frac{285.25}{4 \times 256} = 0.2786\text{ m} = 27.86\text{ cm} \]
৪. দৈর্ঘ্যের প্রয়োজনীয় পরিবর্তন ($\Delta L$):
\[ \Delta L = L_1 - L_2 = 54.02 - 27.86 = 26.16\text{ cm} \]
অতএব, নলের প্রাথমিক দৈর্ঘ্য ছিল $54.02\text{ cm}$ এবং এটিকে বন্ধ নলে রূপান্তর করে ঐকতানে আনতে দৈর্ঘ্য $26.16\text{ cm}$ কমাতে হবে।
\end{tcolorbox}
% ------------------------------------------
% QUESTION SECTION 18 - DETAILED & CLASSY
% ------------------------------------------
\begin{tcolorbox}[questionbox={প্রশ্ন (Question) 6}]
\noindent
একই দিকে ধাবমান দুটি অগ্রগামী শব্দ তরঙ্গের সমীকরণ $y_1 = 0.5 \sin \pi(110.03t - 3.09x)$ এবং $y_2 = 0.5 \sin \pi(134t - 3.4x)$, যেখানে সকল রাশি SI এককে। (ক) তরঙ্গদ্বয়ের উপরিপাতনে সৃষ্ট বিটের পর্যায়কাল কত? (খ) বিট চক্রে সর্বোচ্চ ও সর্বনিম্ন তীব্রতার অনুপাত ($I_{\max} / I_{\min}$) নির্ণয় করো যদি তাদের বিস্তারের অনুপাত $2:1$ হয়।

\vspace{0.8em}
\begin{english}
\noindent\textit{\color{darkgray}
Equations of two progressive sound waves propagating in the same direction are $y_1 = 0.5 \sin \pi(110.03t - 3.09x)$ and $y_2 = 0.5 \sin \pi(134t - 3.4x)$, where all quantities are in SI units. (a) What is the time period of beats formed by their superposition? (b) Find the ratio of maximum to minimum intensity ($I_{\max} / I_{\min}$) in the beat cycle if their amplitudes were in the ratio $2:1$.
}
\end{english}
\end{tcolorbox}

\vspace{1em}

% ------------------------------------------
% SOLUTION SECTION 18
% ------------------------------------------
\begin{tcolorbox}[solutionbox={সমাধান (Solution)}]
\noindent
১. বিট কম্পাঙ্ক ও পর্যায়কাল নির্ণয়: \\
- ১ম তরঙ্গের কৌণিক কম্পাঙ্ক, $\omega_1 = 110.03\pi \implies 2\pi f_1 = 110.03\pi \implies f_1 = 55.015\text{ Hz}$ \\
- ২য় তরঙ্গের কৌণিক কম্পাঙ্ক, $\omega_2 = 134\pi \implies 2\pi f_2 = 134\pi \implies f_2 = 67\text{ Hz}$ \\
প্রতি সেকেন্ডে বিট সংখ্যা, $f_{\text{beat}} = |f_2 - f_1| = 67 - 55.015 = 11.985\text{ Hz} \approx 12\text{ Hz}$। বিট পর্যায়কাল:
\[ T_{\text{beat}} = \frac{1}{f_{\text{beat}}} = \frac{1}{11.985} \approx 0.0834\text{ s} \]
২. সর্বোচ্চ ও সর্বনিম্ন তীব্রতার অনুপাত নির্ণয়: প্রদত্ত বিস্তারের অনুপাত $\frac{A_1}{A_2} = \frac{2}{1}$।
\[ \frac{I_{\max}}{I_{\min}} = \left(\frac{A_1 + A_2}{A_1 - A_2}\right)^2 = \left(\frac{2 + 1}{2 - 1}\right)^2 = 3^2 = 9:1 \]
অতএব, বিটের পর্যায়কাল $0.0834\text{ s}$ এবং সর্বোচ্চ ও সর্বনিম্ন তীব্রতার অনুপাত $9:1$।
\end{tcolorbox}

\vspace{1.5em}

% ------------------------------------------
% QUESTION SECTION 47 - DETAILED & CLASSY
% ------------------------------------------
\begin{tcolorbox}[questionbox={প্রশ্ন (Question) 7}]
\noindent
একটি বিন্দু শব্দ উৎস হতে নির্গত শব্দের ক্ষমতা $P$। উৎস হতে $r_1 = 10\text{ m}$ দূরত্বে শব্দ তীব্রতা স্তর $\beta_1 = 80\text{ dB}$। (প্রমাণ তীব্রতা $I_0 = 10^{-12}\text{ W m}^{-2}$)

\vspace{0.5em}
\noindent
(ক) শব্দ তীব্রতা স্তর ($\beta$) এবং দূরত্বের লগারিদম ($\log_{10} r$)-এর মধ্যে গাণিতিক সম্পর্ক প্রতিপাদন করো। $\log_{10} r = 1$ থেকে $\log_{10} r = 5$ ব্যবধানে $\beta$ বনাম $\log_{10} r$ লেখচিত্রটি অঙ্কন করো এবং এর ঢাল (Slope) নির্ণয় করো।

\vspace{0.3em}
\noindent
(খ) অঙ্কিত লেখচিত্র হতে $r_2 = 100\text{ m}$ (অর্থাৎ $\log_{10} r_2 = 2$) দূরত্বে শব্দ তীব্রতা স্তর বের করো। শব্দ উৎসের ক্ষমতা $P$ হিসাব করো এবং লেখচিত্রে তীব্রতা স্তর শ্রবণসীমায় (Threshold of Hearing) নেমে $0\text{ dB}$ হওয়ার বিন্দুটি চিহ্নিত করে এর দূরত্ব $r_0$ নির্ণয় করো।

\vspace{0.8em}
{\color{darkgray}\itshape
\noindent
A point sound source emits sound with a total power $P$. At a distance $r_1 = 10\text{ m}$, the sound intensity level is observed to be $\beta_1 = 80\text{ dB}$. (Reference intensity $I_0 = 10^{-12}\text{ W m}^{-2}$).

\vspace{0.3em}
\noindent
(a) Derive the mathematical relationship between the sound intensity level ($\beta$) and the logarithm of distance ($\log_{10} r$). Plot the $\beta$ vs $\log_{10} r$ graph in the range $\log_{10} r \in [1, 5]$ and determine its slope.

\vspace{0.3em}
\noindent
(b) From the plotted graph, determine the sound intensity level at $r_2 = 100\text{ m}$ (i.e., $\log_{10} r_2 = 2$). Calculate the sound power $P$ of the source, and identify on the graph the point where the intensity level drops to the threshold of hearing ($0\text{ dB}$) to find the corresponding distance $r_0$.
}
\end{tcolorbox}

% ------------------------------------------
% SOLUTION SECTION 47
% ------------------------------------------
\begin{tcolorbox}[solutionbox={সমাধান (Solution)}]
\noindent
\textbf{(ক) সমীকরণ প্রতিপাদন, লেখচিত্র অঙ্কন ও ঢাল নির্ণয়:}

\vspace{0.5em}
\noindent
বিন্দু উৎস হতে $r$ দূরত্বে শব্দ তীব্রতা:
\[ I(r) = \frac{P}{4\pi r^2} \]

\noindent
শব্দ তীব্রতা স্তরের সমীকরণ:
\[ \beta(r) = 10 \log_{10}\left(\frac{I}{I_0}\right) = 10 \log_{10}\left(\frac{P}{4\pi I_0 r^2}\right) \]
\[ \beta(r) = 10 \log_{10}\left(\frac{P}{4\pi I_0}\right) - 20 \log_{10}(r) \]

\noindent
ধরি, $X = \log_{10}(r)$ এবং ধ্রুবক $C = 10 \log_{10}\left(\frac{P}{4\pi I_0}\right)$। ফলে সমীকরণটি দাঁড়ায়:
\[ \beta = -20 X + C \]
এটি $y = mX + C$ আকারের একটি ঋণাত্মক ঢালবিশিষ্ট সরলরেখা।

\vspace{0.5em}
\noindent
\textbf{লেখচিত্র অঙ্কনের সারণী ($\beta$ vs $\log_{10}r$):}
\begin{itemize}[leftmargin=2em, itemsep=0.2em]
    \item $r_1 = 10\text{ m} \implies X_1 = \log_{10}(10) = 1 \implies \beta_1 = 80\text{ dB}$
    \item $r_2 = 100\text{ m} \implies X_2 = \log_{10}(100) = 2 \implies \beta_2 = 60\text{ dB}$
    \item $r_0 = 10^5\text{ m} \implies X_0 = \log_{10}(10^5) = 5 \implies \beta_0 = 0\text{ dB}$
\end{itemize}

\vspace{0.3em}
\noindent
\textbf{লেখচিত্রের ঢাল (Slope):}
\[ m = \frac{\Delta \beta}{\Delta (\log_{10} r)} = \frac{\beta_2 - \beta_1}{X_2 - X_1} = \frac{60 - 80}{2 - 1} = -20\text{ dB/decade} \]

\vspace{1em}
\noindent
\textbf{(খ) লেখচিত্র হতে মান নির্ণয় ও উৎসের ক্ষমতা $P$ হিসাব:}

\vspace{0.5em}
\noindent
\textbf{১. $r_2 = 100\text{ m}$-এ তীব্রতা স্তর ($\beta_2$):} \\
লেখচিত্রে $X = \log_{10}(100) = 2$ অক্ষীয় বিন্দু বরাবর নির্দেশিত $\beta$-এর মান:
\[ \beta_2 = 60\text{ dB} \]

\noindent
\textbf{২. শব্দ উৎসের ক্ষমতা ($P$):} \\
$r_1 = 10\text{ m}$ এ $\beta_1 = 80\text{ dB}$ ব্যবহার করে:
\[ 80 = 10 \log_{10}\left(\frac{I_1}{10^{-12}}\right) \implies I_1 = 10^{-4}\text{ W m}^{-2} \]
\[ P = I_1 \times 4\pi r_1^2 = 10^{-4} \times 4\pi (10)^2 = 0.04\pi\text{ W} \approx 0.1257\text{ W} \]

\noindent
\textbf{৩. শ্রবণসীমায় ($0\text{ dB}$) দূরত্ব ($r_0$):} \\
লেখচিত্রের $X$-অক্ষ ছেদক বিন্দুতে ($\beta = 0\text{ dB}$):
\[ \log_{10}(r_0) = 5 \implies r_0 = 10^5\text{ m} = 100\text{ km} \]

\vspace{0.8em}
\noindent
\textbf{বিশেষ দ্রষ্টব্য (শব্দ তীব্রতার দুই সীমা):}
\begin{itemize}[leftmargin=2em, itemsep=0.3em]
    \item \textbf{শ্রবণসীমা (Threshold of Hearing):} মানবকান দ্বারা শোনা যায় এমন সর্বনিম্ন শব্দ তীব্রতা হলো $I_0 = 10^{-12}\text{ W m}^{-2}$। এর সাপেক্ষে তীব্রতা স্তর হয় $\beta = 0\text{ dB}$।
    \item \textbf{বেদনাসীমা (Threshold of Pain):} শব্দের তীব্রতা অতিরিক্ত বৃদ্ধি পেয়ে যখন কানে শ্রবণের অনুভূতির পরিবর্তে তীব্র যন্ত্রণার সৃষ্টি করে, তাকে বেদনাসীমা বলে। এর তীব্রতা $I = 1\text{ W m}^{-2}$ এবং তীব্রতা স্তর $\beta = 120\text{ dB}$।
\end{itemize}

\vspace{1.2em}
\begin{center}
\begin{tikzpicture}[scale=1.05, >=stealth]
    % Grid
    \draw[gray!20, thin] (-0.5,-0.5) grid (6.2,4.5);

    % Axes
    \draw[->, thick] (-0.5,0) -- (6.5,0) node[right] {$X = \log_{10}(r)$};
    \draw[->, thick] (0,-0.5) -- (0,4.8) node[above] {$\beta\text{ (dB)}$};

    % Y-Axis Ticks & Labels
    \draw (0.1, 0.8) -- (-0.1, 0.8) node[left, font=\small] {$20$};
    \draw (0.1, 1.6) -- (-0.1, 1.6) node[left, font=\small] {$40$};
    \draw (0.1, 2.4) -- (-0.1, 2.4) node[left, font=\small] {$60$};
    \draw (0.1, 3.2) -- (-0.1, 3.2) node[left, font=\small] {$80$};
    \draw (0.1, 4.0) -- (-0.1, 4.0) node[left, font=\small] {$100$};
    \node[below left, font=\small] at (0,0) {$0$};

    % X-Axis Ticks & Labels
    \foreach \xval in {1, 2, 3, 4, 5} {
        \draw (\xval, 0.1) -- (\xval, -0.1) node[below=2pt, font=\small] {$\xval$};
    }

    % Plotted Line
    \draw[ultra thick, red!80!black] (0, 4.0) -- (5.2, -0.16);

    % Projection Lines
    \draw[dashed, blue!70!black] (0, 3.2) -- (1, 3.2) -- (1, 0);
    \draw[dashed, blue!70!black] (0, 2.4) -- (2, 2.4) -- (2, 0);

    % Points
    \fill[red!80!black] (1, 3.2) circle (2.5pt) node[above right, font=\small] {$(1, 80\text{ dB})$};
    \fill[red!80!black] (2, 2.4) circle (2.5pt) node[above right, font=\small] {$(2, 60\text{ dB})$};
    \fill[teal!80!black] (5, 0) circle (2.5pt) node[above right, font=\small] {$(5, 0\text{ dB}) \implies r_0 = 10^5\text{ m}$};

    % Slope Text
    \node[red!80!black, font=\small\bfseries] at (3.5, 2.2) {ঢাল $m = -20\text{ dB/decade}$};
\end{tikzpicture}
\end{center}
\end{tcolorbox}

% ------------------------------------------
% QUESTION SECTION 32 - DETAILED & CLASSY
% ------------------------------------------
\begin{tcolorbox}[questionbox={প্রশ্ন (Question) 8}]
\noindent
একটি রকেট উড্ডয়ন কেন্দ্র হতে সোজা ওপরের দিকে $150\text{ m s}^{-1}$ সুষম বেগে উঠছিল। রকেটটি থেকে বিস্ফোরক শব্দ উৎপন্ন হওয়ার $5\text{ s}$ এবং $7\text{ s}$ সময় পর দুটি পর্যবেক্ষণ স্টেশন যথাক্রমে শব্দ শুনতে পেল। বিস্ফোরণের ফলে উৎপন্ন শব্দতরঙ্গ $y = 0.1 \sin(69.53\pi t + 0.6283x)$ দ্বারা প্রকাশিত হলে, বিস্ফোরণের মুহূর্তে রকেটের উচ্চতা এবং স্টেশনদ্বয়ের মধ্যবর্তী দূরত্ব নির্ণয় করো।

\vspace{0.8em}
\begin{english}
\noindent\textit{\color{darkgray}
A rocket was ascending vertically upwards at a uniform speed of $150\text{ m s}^{-1}$ from a launchpad. Sound from an explosion on the rocket was received by two observation stations after $5\text{ s}$ and $7\text{ s}$ respectively. If the explosion sound wave is represented by $y = 0.1 \sin(69.53\pi t + 0.6283x)$, calculate the rocket's height at the moment of explosion and the distance between the two stations.
}
\end{english}
\end{tcolorbox}

\vspace{1em}

% ------------------------------------------
% SOLUTION SECTION 32
% ------------------------------------------
\begin{tcolorbox}[solutionbox={সমাধান (Solution)}]
\noindent
১. শব্দের বেগ ($v$) নির্ণয়: \\
প্রদত্ত তরঙ্গ সমীকরণ হতে,
\[ k = 0.6283 \implies \frac{2\pi}{\lambda} = 0.6283 \implies \lambda = 10\text{ m} \]
\[ \omega = 69.53\pi \implies 2\pi f = 69.53\pi \implies f = 34.765\text{ Hz} \]
শব্দের বেগ:
\[ v = f \times \lambda = 34.765 \times 10 = 347.65\text{ m s}^{-1} \approx 348\text{ m s}^{-1} \]

\noindent
২. বিস্ফোরণের মুহূর্তে রকেটের উল্লম্ব উচ্চতা ($h$) নির্ণয়: \\
১ম স্টেশনটি রকেটের ঠিক নিচে অবস্থান করায় এর কাছে শব্দ পৌঁছাতে সর্বনিম্ন সময় ($t_1 = 5\text{ s}$) লাগে। \\
বিস্ফোরণের মুহূর্তে রকেটের উচ্চতা:
\[ h = v \times t_1 = 348 \times 5 = 1740\text{ m} = 1.74\text{ km} \]

\noindent
৩. স্টেশনদ্বয়ের মধ্যবর্তী অনুভূমিক দূরত্ব ($x$) নির্ণয়: \\
২য় স্টেশনে শব্দ পৌঁছাতে সময় লাগে $t_2 = 7\text{ s}$। ২য় স্টেশন হতে বিস্ফোরণস্থলের লম্ব-দূরত্ব (অতিভুজ):
\[ d_2 = v \times t_2 = 348 \times 7 = 2436\text{ m} \]
পিথাগোরাসের উপপাদ্য অনুযায়ী স্টেশনদ্বয়ের অনুভূমিক দূরত্ব:
\[ x = \sqrt{d_2^2 - h^2} = \sqrt{2436^2 - 1740^2} = \sqrt{5934096 - 3027600} = \sqrt{2906496} \approx 1704.84\text{ m} \]

\vspace{0.5em}
\noindent
অতএব, বিস্ফোরণের মুহূর্তে রকেটের উচ্চতা ছিল $1740\text{ m}$ এবং পর্যবেক্ষণ স্টেশনদ্বয়ের মধ্যবর্তী দূরত্ব $1704.84\text{ m}$।
\end{tcolorbox}

% ------------------------------------------
% QUESTION SECTION 48 - DETAILED & CLASSY
% ------------------------------------------
\begin{tcolorbox}[questionbox={প্রশ্ন (Question) 9}]
\noindent
দুটি শব্দ তরঙ্গ $y_1(t) = 0.05 \sin(2\pi f_1 t)$ এবং $y_2(t) = 0.05 \sin(2\pi f_2 t)$ (যেখানে $f_1 > f_2$) উপরিপাতিত হয়ে বিট (Beats) সৃষ্টি করে। উৎপন্ন লব্ধি তরঙ্গের বাহক কম্পাঙ্কের (Carrier frequency) মান $100\text{ Hz}$ এবং প্রতি সেকেন্ডে ৪টি বিট এনভেলপ ডালপালা (Loops) সম্পন্ন হয়।

\vspace{0.5em}
\noindent
(ক) মূল তরঙ্গদ্বয়ের কম্পাঙ্ক $f_1$ ও $f_2$ নির্ণয় করো।

\vspace{0.3em}
\noindent
(খ) বিট এনভেলপের গাণিতিক সমীকরণ ($Y_{\text{envelope}}(t)$) গঠন করো। $t = 0$ থেকে $t = 1\text{ s}$ সময় ব্যবধানে লব্ধি তরঙ্গের সরণ-সময় ($y-t$) ও বিট এনভেলপের লেখচিত্রটি অঙ্কন করো এবং ধ্বংসাত্মক ব্যতিচারের (Destructive interference) ফলে যে মুহূর্তগুলোতে শব্দ তীব্রতা সম্পূর্ণ শূন্য হয়, তা লেখচিত্রে সুনির্দিষ্টভাবে চিহ্নিত করো।

\vspace{0.8em}
{\color{darkgray}\itshape
\noindent
Two sound waves $y_1(t) = 0.05 \sin(2\pi f_1 t)$ and $y_2(t) = 0.05 \sin(2\pi f_2 t)$ (where $f_1 > f_2$) superpose to form beats. The carrier frequency of the resulting wave is observed to be $100\text{ Hz}$ and 4 beat loops are formed per second.

\vspace{0.3em}
\noindent
(a) Determine the frequencies $f_1$ and $f_2$ of the component waves.

\vspace{0.3em}
\noindent
(b) Form the mathematical equation of the beat envelope ($Y_{\text{envelope}}(t)$). Plot the displacement-time ($y-t$) graph of the resultant wave along with the beat envelope in the range $t \in [0, 1\text{ s}]$, and explicitly mark the key points where the sound intensity drops to zero due to destructive interference.
}
\end{tcolorbox}

\vspace{1.2em}

% ------------------------------------------
% SOLUTION SECTION 48
% ------------------------------------------
\begin{tcolorbox}[solutionbox={সমাধান (Solution)}]
\noindent
\textbf{(ক) মূল তরঙ্গদ্বয়ের কম্পাঙ্ক ($f_1, f_2$) নির্ণয়:}

\vspace{0.5em}
\noindent
তরঙ্গের উপরিপাতন নীতি অনুসারে লব্ধি তরঙ্গের সরণ:
\[ y_{\text{net}}(t) = y_1(t) + y_2(t) = 0.05 \left[\sin(2\pi f_1 t) + \sin(2\pi f_2 t)\right] \]
\[ y_{\text{net}}(t) = 2 (0.05) \cos\left(2\pi \frac{f_1 - f_2}{2} t\right) \sin\left(2\pi \frac{f_1 + f_2}{2} t\right) \]

\noindent
প্রদত্ত তথ্যানুসারে, বাহক কম্পাঙ্ক (Carrier frequency):
\[ f_{\text{carrier}} = \frac{f_1 + f_2}{2} = 100\text{ Hz} \implies f_1 + f_2 = 200\text{ Hz} \quad \text{--- (১)} \]

\noindent
প্রতি সেকেন্ডে বিট সংখ্যা (Beat frequency):
\[ f_{\text{beat}} = f_1 - f_2 = 4\text{ Hz} \quad \text{--- (২)} \]

\noindent
সমীকরণ (১) ও (২) যোগ ও বিয়োগ করে পাই:
\[ 2f_1 = 204 \implies f_1 = 102\text{ Hz} \]
\[ 2f_2 = 196 \implies f_2 = 98\text{ Hz} \]

\vspace{1em}
\noindent
\textbf{(খ) বিট এনভেলপের সমীকরণ, শব্দ শূন্য হওয়ার মুহূর্ত ও লেখচিত্র অঙ্কন:}

\vspace{0.5em}
\noindent
লব্ধি তরঙ্গের মডুলেটেড বিস্তার বা এনভেলপের গাণিতিক সমীকরণ:
\[ Y_{\text{envelope}}(t) = \pm 2A \cos\left(2\pi \frac{f_1 - f_2}{2} t\right) \]
\[ Y_{\text{envelope}}(t) = \pm 0.1 \cos\left(2\pi \times \frac{4}{2} t\right) = \pm 0.1 \cos(4\pi t)\text{ m} \]

\vspace{0.3em}
\noindent
ধ্বংসাত্মক ব্যতিচারে শব্দ তীব্রতা সম্পূর্ণ শূন্য হওয়ার জন্য এনভেলপের বিস্তার শূন্য হতে হবে:
\[ \cos(4\pi t) = 0 \implies 4\pi t = (2n + 1)\frac{\pi}{2} \quad [n = 0, 1, 2, 3, \dots] \]
\[ t = \frac{2n + 1}{8}\text{ s} \]

\vspace{0.3em}
\noindent
$t \in [0, 1\text{ s}]$ সীমার মধ্যে লেখচিত্রে চিহ্নিত করার জন্য প্রয়োজনীয় সময়গুলো:
\begin{itemize}[leftmargin=2em, itemsep=0.2em]
    \item $n = 0 \implies t_1 = \frac{1}{8}\text{ s} = 0.125\text{ s}$ \quad (\textbf{শব্দ শূন্য})
    \item $n = 1 \implies t_2 = \frac{3}{8}\text{ s} = 0.375\text{ s}$ \quad (\textbf{শব্দ শূন্য})
    \item $n = 2 \implies t_3 = \frac{5}{8}\text{ s} = 0.625\text{ s}$ \quad (\textbf{শব্দ শূন্য})
    \item $n = 3 \implies t_4 = \frac{7}{8}\text{ s} = 0.875\text{ s}$ \quad (\textbf{শব্দ শূন্য})
\end{itemize}

\vspace{1.2em}
\begin{center}
\begin{tikzpicture}[scale=1.05, >=stealth]
    
    % Grid
    \draw[gray!20, thin] (-0.5,-1.8) grid (6.2,1.8);

    % Axes
    \draw[->, thick] (-0.5,0) -- (6.5,0) node[right] {$t\text{ (s)}$};
    \draw[->, thick] (0,-1.8) -- (0,2.1) node[above] {$y\text{ (m)}$};

    % Y-Axis Ticks & Labels
    \draw (0.1, 1.2) -- (-0.1, 1.2) node[left, font=\small] {$+0.1$};
    \draw (0.1, -1.2) -- (-0.1, -1.2) node[left, font=\small] {$-0.1$};
    \node[below left, font=\small] at (0,0) {$0$};

    % X-Axis Ticks & Labels
    \foreach \xval/\xtext in {0.7/\frac{1}{8}, 1.4/\frac{1}{4}, 2.1/\frac{3}{8}, 2.8/\frac{1}{2}, 3.5/\frac{5}{8}, 4.2/\frac{3}{4}, 4.9/\frac{7}{8}, 5.6/1} {
        \draw (\xval, 0.1) -- (\xval, -0.1) node[below=2pt, font=\small] {$\xtext$};
    }

    % Envelope curves
    \draw[thick, red!80!black, dashed, domain=0:5.6, samples=150] 
        plot (\x, {1.2*cos(900*\x/7)});
    \draw[thick, red!80!black, dashed, domain=0:5.6, samples=150] 
        plot (\x, {-1.2*cos(900*\x/7)});

    % Inside wave (Carrier frequency = 100 Hz)
\draw plot[domain=0:5, samples=500] 
  (\x, {1.2 * cos(128.5714*\x) * sin(360*(\x*17.85714 - floor(\x*17.85714))) });

    % Nodes (Zero intensity points)
    \foreach \xval in {0.7, 2.1, 3.5, 4.9} {
        \fill[teal!80!black] (\xval, 0) circle (2.5pt);
    }

    % Title
    \node[above, blue!80!black, font=\small\bfseries] at (3.8, 2.1) {বিট এনভেলপ ও লব্ধি তরঙ্গ ($y-t$ লেখচিত্র)};
\end{tikzpicture}
\end{center}
\end{tcolorbox}
% ------------------------------------------
% QUESTION SECTION 22 - DETAILED & CLASSY
% ------------------------------------------
\begin{tcolorbox}[questionbox={প্রশ্ন (Question) 10}]
\noindent
গাছের ডালে বসা একটি পাখি অনবরত নির্দিষ্ট শক্তিতে ডাকছে। একজন পর্যবেক্ষক পাখির ডাকের তীব্রতা পর্যবেক্ষণ শুরু করার পর নিজের অবস্থান থেকে পাখিটির দূরত্ব দ্বিগুণ করে নিল। পাখির সৃষ্ট শব্দের ক্ষমতা $10^{-7}\text{ W}$। (ক) পাখির অবস্থান হতে ন্যূনতম কত দূরত্বে শব্দ তীব্রতা স্তর শূন্য ডেসিবল ($0\text{ dB}$) হবে? (খ) দূরত্ব দ্বিগুণ করার ফলে পর্যবেক্ষকের নিকট শব্দ তীব্রতা স্তরের পরিবর্তন কত হবে গাণিতিকভাবে বিশ্লেষণ করো। ($I_0 = 10^{-12}\text{ W m}^{-2}$)

\vspace{0.8em}
\begin{english}
\noindent\textit{\color{darkgray}
A bird sitting on a tree branch chirps continuously with constant power. An observer listening to the chirping doubles his distance from the bird's location. The sound power generated by the bird is $10^{-7}\text{ W}$. (a) At what minimum distance from the source will the sound intensity level be zero decibels ($0\text{ dB}$)? (b) Mathematically analyze the change in sound intensity level at the observer's position due to doubling the distance. ($I_0 = 10^{-12}\text{ W m}^{-2}$)
}
\end{english}
\end{tcolorbox}

\vspace{1em}

% ------------------------------------------
% SOLUTION SECTION 22
% ------------------------------------------
\begin{tcolorbox}[solutionbox={সমাধান (Solution)}]
\noindent
১. $0\text{ dB}$ তীব্রতা স্তরের জন্য দূরত্ব ($r$) নির্ণয়: \\
$0\text{ dB}$ মানে তীব্রতা $I = I_0 = 10^{-12}\text{ W m}^{-2}$। শব্দ ক্ষমতা, $P = 10^{-7}\text{ W}$।
\[ I = \frac{P}{4\pi r^2} \implies 10^{-12} = \frac{10^{-7}}{4\pi r^2} \]
\[ r^2 = \frac{10^{-7}}{4\pi \times 10^{-12}} = \frac{10^5}{4\pi} \approx 7957.7 \implies r \approx 89.2\text{ m} \]
২. দূরত্ব দ্বিগুণ করার ফলে তীব্রতা স্তরের পরিবর্তন ($\Delta \beta$) নির্ণয়: \\
ধরি, প্রাথমিক দূরত্ব $r_1$ এবং পরবর্তী দূরত্ব $r_2 = 2 r_1$। শব্দ তীব্রতার ব্যস্ত বর্গীয় সূত্রানুসারে:
\[ \frac{I_2}{I_1} = \left(\frac{r_1}{r_2}\right)^2 = \left(\frac{r_1}{2r_1}\right)^2 = \frac{1}{4} \]
শব্দ তীব্রতা স্তরের পরিবর্তন:
\[ \Delta \beta = 10 \log_{10} \left(\frac{I_2}{I_1}\right) = 10 \log_{10} \left(\frac{1}{4}\right) = 10 \times (-0.60206) = -6.02\text{ dB} \]
অতএব, ন্যূনতম দূরত্ব $89.2\text{ m}$ এবং দূরত্ব দ্বিগুণ করায় শব্দ তীব্রতা স্তর $6.02\text{ dB}$ হ্রাস পাবে।
\end{tcolorbox}
\end{document}